
\documentclass{article}
\usepackage{amsfonts}

%%%%%%%%%%%%%%%%%%%%%%%%%%%%%%%%%%%%%%%%%%%%%%%%%%%%%%%%%%%%%%%%%%%%%%%%%%%%%%%%%%%%%%%%%%%%%%%%%%%%%%%%%%%%%%%%%%%%%%%%%%%%%%%%%%%%%%%%%%%%%%%%%%%%%%%%%%%%%%%%%%%%%%%%%%%%%%%%%%%%%%%%%%%%%%%%%%%%%%%%%%%%%%%%%%%%%%%%%%%%%%%%%%%
\usepackage{amsmath, amssymb}
\usepackage{geometry}

\setcounter{MaxMatrixCols}{10}
%TCIDATA{OutputFilter=LATEX.DLL}
%TCIDATA{Version=5.50.0.2960}
%TCIDATA{<META NAME="SaveForMode" CONTENT="1">}
%TCIDATA{BibliographyScheme=Manual}
%TCIDATA{LastRevised=Monday, March 10, 2025 07:27:40}
%TCIDATA{<META NAME="GraphicsSave" CONTENT="32">}

\geometry{a4paper, margin=1in}

\input{tcilatex}

\begin{document}

\title{Summary of Today's Discussion on Clifford Algebras and Spinors}
\author{}
\date{}
\maketitle

\section*{Pauli Matrices and $\text{Cl}(3, 0)$}

The Pauli matrices $\sigma_1, \sigma_2, \sigma_3$ generate the real Clifford
algebra $\text{Cl}(3, 0)$, which corresponds to a 3-dimensional Euclidean
space. The Pauli matrices satisfy the anticommutation relations: 
\begin{equation*}
\{ \sigma_i, \sigma_j \} = 2 \delta_{ij} I, 
\end{equation*}
where $\delta_{ij}$ is the Kronecker delta, and $I$ is the identity matrix.
The algebra $\text{Cl}(3, 0)$ is 8-dimensional and has the basis: 
\begin{equation*}
\{ I, \sigma_1, \sigma_2, \sigma_3, \sigma_1 \sigma_2, \sigma_2 \sigma_3,
\sigma_3 \sigma_1, \sigma_1 \sigma_2 \sigma_3 \}. 
\end{equation*}
The even subalgebra of $\text{Cl}(3, 0)$ is isomorphic to the quaternion
algebra $\mathbb{H}$.

\section*{Complexification and $\text{Cl}(2)$}

The complexification of $\text{Cl}(3, 0)$ is: 
\begin{equation*}
\text{Cl}(3, 0) \otimes \mathbb{C} \cong \text{Cl}(3). 
\end{equation*}
In the complex case, $\text{Cl}(3)$ is isomorphic to $\text{Cl}(2)$, the
complex Clifford algebra generated by two elements $f_1, f_2$ satisfying: 
\begin{equation*}
\{ f_i, f_j \} = 2 \delta_{ij} I. 
\end{equation*}
The algebra $\text{Cl}(2)$ is isomorphic to the algebra of $2 \times 2$
complex matrices, $\text{Mat}(2, \mathbb{C})$.

\section*{Connection to Pauli Matrices}

The Pauli matrices $\sigma_1, \sigma_2$ alone generate $\text{Cl}(2)$, since 
$\sigma_3 = i \sigma_1 \sigma_2$ is already contained in the algebra
generated by $\sigma_1$ and $\sigma_2$. Thus, the Pauli matrices, when
considered over the complex numbers, generate the same algebraic structure
as $\text{Cl}(2)$.

\section*{Spin Group and Spinors}

The spin group $\text{Spin}(1, 3)$ acts on spinors in $\mathbb{C}^4$ through
a spinor representation constructed using the Dirac matrices $\gamma_\mu$.
The Dirac matrices satisfy: 
\begin{equation*}
\{ \gamma_\mu, \gamma_\nu \} = 2 \eta_{\mu\nu} I, 
\end{equation*}
where $\eta_{\mu\nu}$ is the Minkowski metric with signature $(+, -, -, -)$.
The generators of $\text{Spin}(1, 3)$ are: 
\begin{equation*}
\Sigma^{\mu\nu} = \frac{1}{4} [\gamma^\mu, \gamma^\nu]. 
\end{equation*}
Under a Lorentz transformation $\Lambda \in \text{SO}(1, 3)$, a spinor $\psi
\in \mathbb{C}^4$ transforms as: 
\begin{equation*}
\psi \mapsto S(\Lambda) \psi, 
\end{equation*}
where $S(\Lambda) \in \text{Spin}(1, 3)$ is a $4 \times 4$ complex matrix.

\section*{Pauli Matrices and $\text{SU}(2)$}

The Pauli matrices represent the Lie algebra $\mathfrak{su}(2)$ of the
special unitary group $\text{SU}(2)$. The generators of $\mathfrak{su}(2)$
are: 
\begin{equation*}
T_i = -\frac{i}{2} \sigma_i, \quad i = 1, 2, 3. 
\end{equation*}
These generators satisfy the commutation relations: 
\begin{equation*}
[T_i, T_j] = \epsilon_{ijk} T_k. 
\end{equation*}
The group $\text{SU}(2)$ acts on 2-component spinors $\psi \in \mathbb{C}^2$
as: 
\begin{equation*}
\psi \mapsto U(\theta, \mathbf{n}) \psi, 
\end{equation*}
where $U(\theta, \mathbf{n}) = \exp\left( -i \frac{\theta}{2} \mathbf{n}
\cdot \boldsymbol{\sigma} \right)$ represents a rotation by an angle $\theta$
around the axis $\mathbf{n}$.

\section*{Summary}

- The Pauli matrices generate the real Clifford algebra $\text{Cl}(3,0)$,
whose even subalgebra is isomorphic to the quaternion algebra $\mathbb{H}$. 

- The complexification of $\text{Cl}(3,0)$ is isomorphic to $\text{Cl}(2)$,
the complex Clifford algebra generated by two elements. 

- The Pauli matrices, when considered over the complex numbers, generate $%
\text{Cl}(2)$, which is isomorphic to $\text{Mat}(2,\mathbb{C})$. 

- The spin group $\text{Spin}(1,3)$ acts on 4-component Dirac spinors in $%
\mathbb{C}^{4}$ through the Dirac matrices. 

- The Pauli matrices also represent the Lie algebra $\mathfrak{su}(2)$ of $%
\text{SU}(2)$, which acts on 2-component spinors in $\mathbb{C}^{2}$.

\end{document}
